% 1. Настройка формата 16:9
\documentclass[10pt, aspectratio=169, unicode]{beamer}

% --- Шрифты и Язык (Специфика XeLaTeX) ---
\usepackage{fontspec}
\usepackage{xunicode}
\usepackage{xltxtra}
\usepackage{polyglossia}
\setdefaultlanguage{russian}
\setotherlanguage{english}

% --- НАСТРОЙКА ШРИФТОВ (Исправленная версия) ---
\setmainfont{Times New Roman} 
\setsansfont{Arial} 
\setmonofont{Courier New}

% Добавляем явное указание шрифтов для кириллицы (решает твою ошибку)
\newfontfamily\russianfont{Times New Roman}
\newfontfamily\russianfontsf{Arial}
\newfontfamily\russianfonttt{Courier New} 

\usefonttheme{serif}

% --- Цвета и Тема (Стиль МАИ) ---
% Используем тему Madrid, но меняем цвета
\usetheme{Madrid}

% Определяем цвет МАИ (примерный голубой)
\definecolor{maiblue}{RGB}{0, 114, 188} 
\definecolor{maigray}{RGB}{240, 240, 240}

% Применяем цвета
\usecolortheme[named=maiblue]{structure} % Основной цвет элементов
\setbeamercolor{frametitle}{bg=maiblue, fg=white} % Заголовок слайда
\setbeamercolor{title in head/foot}{bg=maiblue, fg=white} % Нижний бар

% Убираем навигационные значки (они обычно мешают)
\setbeamertemplate{navigation symbols}{}

% --- Пакеты ---
\usepackage{graphicx}  % Картинки
\usepackage{booktabs}  % Красивые таблицы
\usepackage{tabularx}  % ← ДОБАВИТЬ — таблицы с X-колонками (слайд 4)
\usepackage{amsmath}   % ← ДОБАВИТЬ — формулы (слайд 7)
\usepackage{caption}   % Настройка подписейep=endash, name=Рис.}
\captionsetup[table]{labelformat=simple, labelsep=endash, name=Табл.}

% --- Оформление кода (Python) ---
\usepackage{listings}
\usepackage{xcolor}

% Цвета для кода
\definecolor{codegreen}{rgb}{0,0.6,0}
\definecolor{codegray}{rgb}{0.5,0.5,0.5}
\definecolor{codepurple}{rgb}{0.58,0,0.82}
\definecolor{backcolour}{rgb}{0.95,0.95,0.92}

\lstdefinestyle{pythonstyle}{
    backgroundcolor=\color{backcolour},   
    commentstyle=\color{codegreen},
    keywordstyle=\color{magenta},
    numberstyle=\tiny\color{codegray},
    stringstyle=\color{codepurple},
    basicstyle=\ttfamily\footnotesize, % Шрифт кода
    breakatwhitespace=false,         
    breaklines=true,                 
    captionpos=b,                    
    keepspaces=true,                 
    numbers=left,                    % Нумерация строк слева
    numbersep=5pt,                  
    showspaces=false,                
    showstringspaces=false,
    showtabs=false,                  
    tabsize=2,
    language=Python                  % Язык по умолчанию
}
\lstset{
  extendedchars=true,
  keepspaces=true,
  % Этот блок сопоставляет символы, чтобы xelatex их не терял
  literate={а}{{\cyra}}1 {б}{{\cyrb}}1 {в}{{\cyrv}}1 {г}{{\cyrg}}1 {д}{{\cyrd}}1 {е}{{\cyre}}1 {ё}{{\"{\cyre}}}1 {ж}{{\cyrzh}}1 {з}{{\cyrz}}1 {и}{{\cyri}}1 {й}{{\cyrj}}1 {к}{{\cyrk}}1 {л}{{\cyrl}}1 {м}{{\cyrm}}1 {н}{{\cyrn}}1 {о}{{\cyro}}1 {п}{{\cyrp}}1 {р}{{\cyrr}}1 {с}{{\cyrs}}1 {т}{{\cyrt}}1 {у}{{\cyru}}1 {ф}{{\cyrf}}1 {х}{{\cyrh}}1 {ц}{{\cyrc}}1 {ч}{{\cyrch}}1 {ш}{{\cyrsh}}1 {щ}{{\cyrshch}}1 {ъ}{{\cyrhrdsn}}1 {ы}{{\cyry}}1 {ь}{{\cyrsftsn}}1 {э}{{\cyre}}1 {ю}{{\cyryu}}1 {я}{{\cyrya}}1
           {А}{{\CYRA}}1 {Б}{{\CYRB}}1 {В}{{\CYRV}}1 {Г}{{\CYRG}}1 {Д}{{\CYRD}}1 {Е}{{\CYRE}}1 {Ё}{{\"{\CYRE}}}1 {Ж}{{\CYRZH}}1 {З}{{\CYRZ}}1 {И}{{\CYRI}}1 {Й}{{\CYRJ}}1 {К}{{\CYRK}}1 {Л}{{\CYRL}}1 {М}{{\CYRM}}1 {Н}{{\CYRN}}1 {О}{{\CYRO}}1 {П}{{\CYRP}}1 {Р}{{\CYRR}}1 {С}{{\CYRS}}1 {Т}{{\CYRT}}1 {У}{{\CYRU}}1 {Ф}{{\CYRF}}1 {Х}{{\CYRH}}1 {Ц}{{\CYRC}}1 {Ч}{{\CYRCH}}1 {Ш}{{\CYRSH}}1 {Щ}{{\CYRSHCH}}1 {Ъ}{{\CYRHRDSN}}1 {Ы}{{\CYRY}}1 {Ь}{{\CYRSFTSN}}1 {Э}{{\CYRE}}1 {Ю}{{\CYRYU}}1 {Я}{{\CYRYA}}1
}


% ─────────────────────────────────────────────
% --- Данные титульного листа ---
\title[ВКР]{Разработка интеллектуального модуля интерактивного анализа данных в реляционных базах данных}
\subtitle{Выпускная квалификационная работа}
\author[Гимазетдинов Д.Р.]{Студент: Гимазетдинов Дмитрий Русланович}
\institute[МАИ]{Московский авиационный институт\\ (Национальный исследовательский университет)}
\date{2026}
% ─────────────────────────────────────────────

\begin{document}

% ══════════════════════════════════════════════
% 1. Титульный лист
% ══════════════════════════════════════════════
\begin{frame}
  \titlepage
\end{frame}

% ══════════════════════════════════════════════
% 2. Актуальность темы
% ══════════════════════════════════════════════
\begin{frame}{Актуальность темы}
  \begin{columns}[T]
    \begin{column}{0.55\textwidth}
      \textbf{Проблема:}
      \begin{itemize}
        \item Бизнес-цели требуют быстрого отклика от аналитиков
        \item Аналитики вынуждены ждать разработчиков для получения отчётов
        \item Прослойка между аналитиком и данными становится рудиментом
      \end{itemize}
      \vspace{6pt}
      \textbf{Пример:}\\
      {\small«Какие самолёты чаще всего летали из Москвы в Санкт-Петербург в прошлом месяце с загрузкой более 80\%?» — требует нескольких JOIN и знания схемы}
    \end{column}
    \begin{column}{0.42\textwidth}
      \textbf{Факторы актуальности:}
      \begin{enumerate}
        \item Экспоненциальный рост объёмов данных
        \item Высокий порог входа в аналитику (SQL)
        \item Активное развитие LLM-моделей
        \item Ограничения существующих решений
      \end{enumerate}
      \vspace{6pt}
      \textbf{Решение:} NLIDB-интерфейс на основе RAG + LLM + агентный подход + режим самокоррекции
    \end{column}
  \end{columns}
\end{frame}

% ══════════════════════════════════════════════
% 3. Цель и задачи
% ══════════════════════════════════════════════
\begin{frame}{Цель и задачи исследования}
  \textbf{Цель:}\\[4pt]
  Разработать программное решение, которое устранит барьер между пользователями без технического образования и сложными реляционными базами данных.

  \vspace{8pt}
  \textbf{Задачи:}
  \begin{enumerate}
    \item Провести оценку современных методов трансляции естественного языка в SQL
    \item Спроектировать RAG-архитектуру для работы со схемой БД (метаинформация, связи таблиц)
    \item Разработать механизм самокоррекции ИИ-агента (Self-Heal)
    \item Провести оценку качества генерации до и после самокоррекции
    \item Сравнить LLM-провайдеры: \texttt{GLM-4.7}, \texttt{Qwen3-Coder-Next}, \texttt{GigaChat-2}
    \item Реализовать микросервисную архитектуру системы
  \end{enumerate}
\end{frame}

% ══════════════════════════════════════════════
% 4. Сравнение существующих решений
% ══════════════════════════════════════════════
\begin{frame}{Сравнение с существующими решениями}
  \begin{table}
    \centering
    \small
    \begin{tabularx}{\textwidth}{|>{\raggedright\arraybackslash}X|c|c|c|c|}
      \hline
      \textbf{Функционал} & \textbf{Vanna.AI} & \textbf{LangChain} & \textbf{Metabase} & \textbf{Наш продукт} \\
      \hline
      Генерация SQL (Text-to-SQL)      & + & + & + & \textbf{+} \\
      \hline
      RAG для схемы БД                 & + & Ограничено & -- & \textbf{+} \\
      \hline
      Локальный запуск (Privacy)       & + & + & -- & \textbf{+} \\
      \hline
      Механизм самокоррекции (Self-Heal) & -- & -- & -- & \textbf{+} \\
      \hline
      Авто-визуализация данных         & $\pm$ & -- & + & В разработке \\
      \hline
    \end{tabularx}
  \end{table}

  \vspace{6pt}
  \begin{block}{Ключевые преимущества нашего решения}
    \begin{itemize}
      \item \textbf{Приватность:} можно все вычисления вести локально, данные не покидают контур
      \item \textbf{Self-Heal:} автоматическая итеративная самокоррекция SQL
      \item \textbf{RAG для сложных схем:} оптимизировано для многотабличных БД
    \end{itemize}
  \end{block}
\end{frame}

% ══════════════════════════════════════════════
% 5.2 Архитектура системы
% ══════════════════════════════════════════════
\begin{frame}{Архитектура системы}
  \begin{columns}[T]
    \begin{column}{0.48\textwidth}
      \textbf{Микросервисы:}
      \begin{itemize}
        \item \textbf{Frontend} — веб-интерфейс пользователя
        \item \textbf{RAG Service} — оркестратор пайплайна
        \item \textbf{Agent Module} — LangGraph агент (Chain-of-Thought)
        \item \textbf{Embedding Model Service} — локальная векторизация (\texttt{multilingual-e5-base})
        \item \textbf{Retrieval Pipeline} — семантический поиск в ChromaDB
        \item \textbf{Ingest Pipeline} — индексация документов (Redis Queue)
        \item \textbf{DB Connector Service} — безопасный доступ к PostgreSQL
      \end{itemize}
    \end{column}
    \begin{column}{0.48\textwidth}
      \textbf{Хранилища:}
      \begin{itemize}
        \item \textbf{MinIO} — объектное хранилище файлов (S3-совместимо)
        \item \textbf{ChromaDB} — векторная БД
        \item \textbf{PostgreSQL} — реляционная БД
        \item \textbf{Redis} — очередь задач
      \end{itemize}
    \end{column}
  \end{columns}
\end{frame}

% ══════════════════════════════════════════════
% 5.2. Архитектура системы
% ══════════════════════════════════════════════
\begin{frame}{Архитектура системы}

    \begin{figure}[H]
        \centering
        \includegraphics[width=0.65\textwidth]{/Users/cube/Documents/Development/Наработки по дипломной работе/Диплом/tex/prezentation/materials/HLD.drawio.png}
        \caption{HLD диаграмма архитектуры системы}
        \label{hld}
    \end{figure}

\end{frame}

% ══════════════════════════════════════════════
% 5.3. Архитектура системы
% ══════════════════════════════════════════════
\begin{frame}{Архитектура системы}

    \begin{figure}[H]
        \centering
        \includegraphics[width=0.7\textwidth]{/Users/cube/Documents/Development/Наработки по дипломной работе/Диплом/tex/prezentation/materials/rag_pipeline.drawio.png}
        \caption{RAG-пайплайн}
        \label{rag}
    \end{figure}

\end{frame}

% ══════════════════════════════════════════════
% 6.1 Механизм самокоррекции
% ══════════════════════════════════════════════
\begin{frame}{Механизм самокоррекции}
  \begin{columns}[T]
    \begin{column}{0.44\textwidth}
      \textbf{Агентный цикл (LangGraph):}
      \begin{enumerate}
        \item Переписывание запроса для RAG-поиска
        \item Семантический поиск контекста схемы БД
        \item Генерация SQL на основе контекста
        \item Выполнение SQL в DB Connector Service
        \item \textbf{Критика результата} — LLM анализирует ошибку и выносит вердикт
        \item Если нужно: запрос схемы проблемных таблиц
        \item Улучшение SQL с учётом замечаний
        \item Повтор до 5 итераций
      \end{enumerate}
    \end{column}
    \begin{column}{0.55\textwidth}
        \begin{figure}[H]
            \centering
            \includegraphics[width=0.55\textwidth]{/Users/cube/Documents/Development/Наработки по дипломной работе/Диплом/tex/prezentation/materials/self_correction.png}
            \caption{Алгоритм самокоррекции}
            \label{self_correction}
        \end{figure}
    \end{column}
  \end{columns}
\end{frame}

% ══════════════════════════════════════════════
% 7. Метрики сравнения LLM
% ══════════════════════════════════════════════
\begin{frame}{Метрики оценки качества генерации SQL}
  \begin{columns}[T]
    \begin{column}{0.48\textwidth}
      \textbf{Execution Accuracy (EA):}
      \[
        \text{EA} = \frac{|\{q : \text{exec}(q) = \text{exec}(q^*)\}|}{|Q|}
      \]
      Доля запросов, результат выполнения которых совпал с эталоном. Совпадение результирующих множеств.

      \vspace{8pt}
      \textbf{Exact Matching (EM):}
      \[
        \text{EM} = \frac{|\{q : q = q^*\}|}{|Q|}
      \]
      Точное совпадение SQL текста.
    \end{column}
    \begin{column}{0.48\textwidth}
      \textbf{Exact Set Matching (ESM):}
      \[
        \text{ESM} = \frac{|\{q : R(q) \equiv R(q^*)\}|}{|Q|}
      \]
      Cравнивает структурные компоненты SQL-запроса (например, списки столбцов, условия фильтрации, ...)

      \vspace{8pt}
      \textbf{Intellect (IEA):}
      \[
        \text{IEA} = \frac{|\{q : \text{exec}(q) \approx \text{exec}(q^*)\}|}{|Q|}
      \]
      Учитывает семантически верные, но синтаксически отличающиеся запросы.

      \vspace{4pt}
      {\small Тестовый набор: \textbf{31 SQL-задача} разного уровня сложности на БД «Авиаперевозки»}
    \end{column}
  \end{columns}
\end{frame}

% ══════════════════════════════════════════════
% 8.1 Графики метрик
% ══════════════════════════════════════════════
\begin{frame}{Результаты сравнения LLM-моделей}
    \begin{minipage}{0.49\textwidth}
        \centering
        \includegraphics[width=\textwidth]{/Users/cube/Documents/Development/Наработки по дипломной работе/Диплом/tex/prezentation/materials/llm_before.png}
        % \captionof{figure}{До самокоррекции}
    \end{minipage}\hfill
    \begin{minipage}{0.49\textwidth}
        \centering
        \includegraphics[width=\textwidth]{/Users/cube/Documents/Development/Наработки по дипломной работе/Диплом/tex/prezentation/materials/llm_after.png}
        % \captionof{figure}{После самокоррекции}
    \end{minipage}
\end{frame}

% ══════════════════════════════════════════════
% 8.2 Графики метрик
% ══════════════════════════════════════════════
\begin{frame}{Результаты сравнения LLM-моделей}
  \begin{columns}[T]
    \begin{column}{0.48\textwidth}
      \centering
      \textbf{Без самокоррекции}
      \begin{table}
        \small
        \begin{tabular}{|l|c|c|c|c|}
          \hline
          \textbf{Модель} & \textbf{EA} & \textbf{EM} & \textbf{ESM} & \textbf{IEA} \\
          \hline
          GLM-4.7      & \textbf{54,8} & \textbf{25,8} & \textbf{35,5} & \textbf{64,5} \\
          Qwen3-Coder  & 45,2 & 19,4 & 29,0 & 54,8 \\
          GigaChat-2   & 16,1 &  6,5 &  9,7 & 22,6 \\
          \hline
        \end{tabular}
      \end{table}
    \end{column}
    \begin{column}{0.48\textwidth}
      \centering
      \textbf{С самокоррекцией}
      \begin{table}
        \small
        \begin{tabular}{|l|c|c|c|c|}
          \hline
          \textbf{Модель} & \textbf{EA} & \textbf{EM} & \textbf{ESM} & \textbf{IEA} \\
          \hline
          GLM-4.7      & \textbf{67,7} & \textbf{38,7} & \textbf{48,4} & \textbf{77,4} \\
          Qwen3-Coder  & 61,3 & 32,3 & 41,9 & 71,0 \\
          GigaChat-2   & 29,0 & 12,9 & 22,6 & 38,7 \\
          \hline
        \end{tabular}
      \end{table}
    \end{column}
  \end{columns}

  \vspace{6pt}
  \begin{block}{Прирост EA от самокоррекции}
    \centering
    \begin{tabular}{lccc}
      \toprule
      & GLM-4.7 & Qwen3-Coder & GigaChat-2 \\
      \midrule
      $\Delta$ EA & \textbf{+13,0 п.п.} & +16,1 п.п. & +12,9 п.п. \\
      \bottomrule
    \end{tabular}
  \end{block}

  {\small \textbf{Вывод:} GLM-4.7 лидирует по всем метрикам в обоих режимах. Qwen демонстрирует наибольший относительный прирост от самокоррекции. GigaChat-2 существенно уступает.}
\end{frame}

% ══════════════════════════════════════════════
% 9. Результаты и внедрение
% ══════════════════════════════════════════════
\begin{frame}{Результаты и внедрение}
  \begin{columns}[T]
    \begin{column}{0.48\textwidth}
      \textbf{Технические характеристики:}
      \begin{itemize}
        \item Векторизация запроса: \textbf{80–150 мс}
        \item Семантический поиск (10K чанков): \textbf{50–100 мс}
        \item Полный агентный цикл (1–3 итерации): \textbf{15–45 сек}
        \item Индексация документа (10 стр.): \textbf{20–40 сек}
      \end{itemize}

      \vspace{6pt}
      \textbf{Развёртывание:}
      \begin{itemize}
        \item Одна команда: \texttt{docker compose up -d}
        \item Веб-интерфейс: \texttt{localhost:3000}
        \item API: \texttt{localhost:8003}
        \item Поддерживаемые форматы: PDF, DOCX, TXT
      \end{itemize}
    \end{column}
    \begin{column}{0.48\textwidth}
      \textbf{Область применения:}
      \begin{itemize}
        \item Розничная торговля и логистика
        \item Финансовый сектор
        \item Производство и телеком
        \item Авиаперевозки
        \item Банковские хранилища
      \end{itemize}

      \vspace{6pt}
      \textbf{Апробация:}\\
      Тестирование на демонстрационной БД «Авиаперевозки» (Postgres Pro):
      \begin{itemize}
        \item 8 взаимосвязанных таблиц
        \item JSONB-поля с мультиязычными атрибутами
        \item 31 SQL-задача разного уровня сложности
      \end{itemize}
    \end{column}
  \end{columns}
\end{frame}

% ══════════════════════════════════════════════
% 10. Практическая значимость
% ══════════════════════════════════════════════
\begin{frame}{Практическая значимость}
  \begin{columns}[T]
    \begin{column}{0.48\textwidth}
      \textbf{Научная новизна:}
      \begin{itemize}
        \item Адаптивный метод семантического поиска контекста схемы БД с многоуровневым индексированием (структура + комментарии + примеры данных)
        \item Оригинальный алгоритм многошаговой самокоррекции SQL («критик» + «исполнитель»)
      \end{itemize}

      \vspace{6pt}
      \textbf{Конфиденциальность данных:}
      \begin{itemize}
        \item Все вычисления могут исполняться локально, внутри корпоративного контура
        \item Модели эмбеддингов запускаются на собственном сервере
        \item Данные не передаются во внешние облачные сервисы
      \end{itemize}
    \end{column}
    \begin{column}{0.48\textwidth}
      \textbf{Масштабируемость:}
      \begin{itemize}
        \item Горизонтальное масштабирование каждого сервиса независимо
        \item ChromaDB: от десятков тысяч до миллионов векторов
        \item Замена LLM или модели эмбеддингов без переработки системы
      \end{itemize}
    \end{column}
  \end{columns}
\end{frame}

% ══════════════════════════════════════════════
% 11. Пример работы системы (заглушка)
% ══════════════════════════════════════════════
\begin{frame}{Пример работы системы}

    \begin{figure}[H]
        \centering
        \includegraphics[width=0.66\textwidth]{/Users/cube/Documents/Development/Наработки по дипломной работе/Диплом/tex/prezentation/materials/preview.png}
        \caption{UI}
        \label{ui}
    \end{figure}

\end{frame}

% ══════════════════════════════════════════════
% 12. Ссылка на GitHub (заглушка)
% ══════════════════════════════════════════════
\begin{frame}{Исходный код проекта}
    \begin{columns}[T]
        \begin{column}{0.48\textwidth}
            \begin{figure}[H]
                \centering
                
\includegraphics[width=0.7\textwidth]{/Users/cube/Documents/Development/Наработки по дипломной работе/Диплом/tex/prezentation/materials/QR-for-C4be.png}
                \caption{GitHub}
                \label{qr1}
            \end{figure}
        \end{column}
        \begin{column}{0.48\textwidth}
            \begin{figure}[H]
                \centering
                
\includegraphics[width=0.7\textwidth]{/Users/cube/Documents/Development/Наработки по дипломной работе/Диплом/tex/prezentation/materials/gitverse.png}
                \caption{GitVerse}
                \label{qr1}
            \end{figure}
        \end{column}
  \end{columns}
   

\end{frame}

\end{document}